% Source: em_algorithm.tex
\begin{notechapter}
% ---- chapter-local definitions ----
\definecolor{codegreen}{rgb}{0,0.6,0}
\definecolor{codegray}{rgb}{0.5,0.5,0.5}
\definecolor{codepurple}{rgb}{0.58,0,0.82}
\definecolor{backcolour}{rgb}{0.95,0.95,0.92}
\lstdefinestyle{mystyle}{
    backgroundcolor=\color{backcolour},   
    commentstyle=\color{codegreen},
    keywordstyle=\color{magenta},
    numberstyle=\tiny\color{codegray},
    stringstyle=\color{codepurple},
    basicstyle=\ttfamily\footnotesize,
    breakatwhitespace=false,         
    breaklines=true,                 
    captionpos=b,                    
    keepspaces=true,                 
    numbers=left,                    
    numbersep=5pt,                  
    showspaces=false,                
    showstringspaces=false,
    showtabs=false,                  
    tabsize=2
}
\lstset{style=mystyle}
% -----------------------------------
\chapter{The EM Algorithm}\label{ch:em-algorithm}
\chaptersubtitle{The three-coin model, from first principles to implementation}

\section*{Introduction for Students}

Welcome! This document will guide you through one of the most beautiful ideas in machine learning: the Expectation-Maximization (EM) algorithm. We'll use a simple three-coin example to understand how EM helps us estimate parameters when we have hidden variables. Even if you're new to these concepts, we'll build everything step by step.

\section{The Three-Coin Problem: Setting the Stage}

\subsection{The Physical Setup}

Imagine you have three coins:
\begin{itemize}
    \item \textbf{Coin A}: Probability of heads = $\pi$ (unknown)
    \item \textbf{Coin B}: Probability of heads = $p$ (unknown)
    \item \textbf{Coin C}: Probability of heads = $q$ (unknown)
\end{itemize}

Here's the experiment:
\begin{enumerate}
    \item Toss coin A
    \item If coin A shows heads, select coin B; if tails, select coin C
    \item Toss the selected coin (B or C) and record the result
    \item Repeat 10 times
\end{enumerate}

You observe these results: 1, 1, 0, 1, 0, 0, 1, 0, 1, 1 (1 = heads, 0 = tails)

\subsection{The Mathematical Challenge}

Here's the tricky part: \textbf{You only see the final coin toss results, not which coin (B or C) was used!}

Let's define our variables mathematically:
\begin{itemize}
    \item $y$: observed outcome (1 = heads, 0 = tails)
    \item $z$: hidden/latent variable (1 = coin B was used, 0 = coin C was used)
    \item $\theta = (\pi, p, q)$: parameters we want to estimate
\end{itemize}

\section{Understanding the Probabilities}

\subsection{Complete Data Likelihood}

If we could see both $y$ and $z$, the probability would be straightforward:

\begin{align*}
P(y, z|\theta) &= P(z|\theta) \cdot P(y|z,\theta) \\
&= \begin{cases}
\pi \cdot p^y(1-p)^{1-y} & \text{if } z = 1 \text{ (coin B was used)} \\
(1-\pi) \cdot q^y(1-q)^{1-y} & \text{if } z = 0 \text{ (coin C was used)}
\end{cases}
\end{align*}

We can write this more compactly using exponents:
\[
P(y, z|\theta) = [\pi p^y(1-p)^{1-y}]^z \cdot [(1-\pi) q^y(1-q)^{1-y}]^{1-z}
\]

\textbf{Why this works}: When $z=1$, the first term is active and the second term becomes $[\cdots]^0 = 1$. When $z=0$, the opposite happens.

\subsection{Marginal Likelihood (What We Actually See)}

Since we don't observe $z$, we need to \textbf{marginalize} over it (sum over all possible values):

\begin{align*}
P(y|\theta) &= \sum_{z \in \{0,1\}} P(y, z|\theta) \\
&= P(y, z=1|\theta) + P(y, z=0|\theta) \\
&= \pi p^y(1-p)^{1-y} + (1-\pi) q^y(1-q)^{1-y}
\end{align*}

For multiple observations $Y = (y_1, \ldots, y_n)$, since they're independent:
\[
P(Y|\theta) = \prod_{j=1}^n \left[\pi p^{y_j}(1-p)^{1-y_j} + (1-\pi) q^{y_j}(1-q)^{1-y_j}\right]
\]

\section{The Maximum Likelihood Estimation Problem}

\subsection{The Log-Likelihood Function}

We want to find parameters that make our observations most likely:
\[
\hat{\theta} = \arg\max_{\theta} \log P(Y|\theta)
\]

The log-likelihood is:
\[
\mathcal{L}(\theta) = \sum_{j=1}^n \log\left[\pi p^{y_j}(1-p)^{1-y_j} + (1-\pi) q^{y_j}(1-q)^{1-y_j}\right]
\]

\subsection{Why This is Hard to Optimize}

Look at the term inside the log: $\pi p^{y_j}(1-p)^{1-y_j} + (1-\pi) q^{y_j}(1-q)^{1-y_j}$

\begin{itemize}
    \item It's a \textbf{sum} inside a \textbf{logarithm}
    \item The parameters $\pi, p, q$ are all mixed together
    \item Taking derivatives leads to messy equations
    \item There's no closed-form solution
\end{itemize}

This is exactly the kind of problem where the EM algorithm shines!

\section{The Brilliant Idea: Expectation-Maximization}

\subsection{Key Insight: Work with Complete Data}

If we knew the hidden variables $z$, the problem would be easy! The complete data log-likelihood would be:
\[
\log P(Y, Z|\theta) = \sum_{j=1}^n \left[z_j \log(\pi p^{y_j}(1-p)^{1-y_j}) + (1-z_j) \log((1-\pi) q^{y_j}(1-q)^{1-y_j})\right]
\]

This separates nicely and we could estimate parameters easily.

\subsection{The Expectation Trick}

Since we don't know $z$, let's take expectations! For any probability distribution $q(z)$, we can write:

\begin{align*}
P(y|\theta) &= \sum_z P(y,z|\theta) \quad \text{(definition)} \\
&= \sum_z q(z) \cdot \frac{P(y,z|\theta)}{q(z)} \quad \text{(multiply and divide by $q(z)$)} \\
&= \mathbb{E}_{q(z)}\left[\frac{P(y,z|\theta)}{q(z)}\right] \quad \text{(this is the key insight!)}
\end{align*}

\textbf{Let's verify this works}:
\begin{align*}
\mathbb{E}_{q(z)}\left[\frac{P(y,z|\theta)}{q(z)}\right] 
&= q(z=1)\frac{P(y,z=1|\theta)}{q(z=1)} + q(z=0)\frac{P(y,z=0|\theta)}{q(z=0)} \\
&= P(y,z=1|\theta) + P(y,z=0|\theta) \\
&= P(y|\theta) \quad \text{✓}
\end{align*}

The $q(z)$ terms cancel out! This works for \textbf{any} distribution $q(z)$.

\subsection{Jensen's Inequality and the Lower Bound}

Now take logarithms of both sides:
\[
\log P(y|\theta) = \log \mathbb{E}_{q(z)}\left[\frac{P(y,z|\theta)}{q(z)}\right]
\]

Since $\log$ is a \textbf{concave} function, Jensen's inequality tells us:
\[
\log \mathbb{E}[X] \geq \mathbb{E}[\log X]
\]

Applying this:
\[
\log P(y|\theta) \geq \mathbb{E}_{q(z)}\left[\log\frac{P(y,z|\theta)}{q(z)}\right]
\]

Let's expand the right-hand side:
\begin{align*}
\mathbb{E}_{q(z)}\left[\log\frac{P(y,z|\theta)}{q(z)}\right] 
&= \mathbb{E}_{q(z)}\left[\log P(y,z|\theta) - \log q(z)\right] \\
&= \mathbb{E}_{q(z)}\left[\log P(y,z|\theta)\right] - \mathbb{E}_{q(z)}\left[\log q(z)\right]
\end{align*}

Define:
\begin{align*}
Q(\theta, q) &= \mathbb{E}_{q(z)}\left[\log P(y,z|\theta)\right] \quad \text{(expected complete data log-likelihood)} \\
H(q) &= -\mathbb{E}_{q(z)}\left[\log q(z)\right] \quad \text{(entropy of $q$)}
\end{align*}

So we have:
\[
\log P(y|\theta) \geq Q(\theta, q) + H(q)
\]

\subsection{The Evidence Lower Bound (ELBO)}

The quantity $\text{ELBO}(\theta, q) = Q(\theta, q) + H(q)$ is called the \textbf{Evidence Lower Bound}. The name makes sense:
\begin{itemize}
    \item \textbf{Evidence}: $\log P(y|\theta)$ is often called the log-evidence
    \item \textbf{Lower Bound}: Our expression is always ≤ the true log-evidence
    \item \textbf{Bound}: It bounds the true likelihood from below
\end{itemize}

\subsection{Finding the Tightest Bound}

When is the bound tight? Jensen's inequality becomes equality when:
\[
\frac{P(y,z|\theta)}{q(z)} = \text{constant} \quad \text{(doesn't depend on $z$)}
\]

If $\frac{P(y,z|\theta)}{q(z)} = c$, then:
\[
P(y,z|\theta) = c \cdot q(z)
\]

Sum over $z$:
\[
\sum_z P(y,z|\theta) = c \cdot \sum_z q(z) \Rightarrow P(y|\theta) = c \cdot 1 \Rightarrow c = P(y|\theta)
\]

So:
\[
P(y,z|\theta) = P(y|\theta) \cdot q(z) \Rightarrow q(z) = \frac{P(y,z|\theta)}{P(y|\theta)} = P(z|y,\theta)
\]

\textbf{The optimal choice is the posterior distribution!}

\section{The EM Algorithm for Three Coins}

\subsection{The EM Strategy}

Since we don't know the optimal $\theta$, we use an iterative approach:

\begin{enumerate}
    \item Start with initial guess $\theta^{(0)}$
    \item Repeat until convergence:
    \begin{itemize}
        \item \textbf{E-step}: Given current $\theta^{(t)}$, compute $q(z) = P(z|y,\theta^{(t)})$
        \item \textbf{M-step}: Find $\theta^{(t+1)} = \arg\max_{\theta} Q(\theta, \theta^{(t)})$
    \end{itemize}
\end{enumerate}

\subsection{E-step: Computing Responsibilities}

For each observation $y_j$, compute:
\[
\gamma_j^{(t)} = P(z_j = 1|y_j, \theta^{(t)}) = \frac{\pi^{(t)} (p^{(t)})^{y_j}(1-p^{(t)})^{1-y_j}}{\pi^{(t)} (p^{(t)})^{y_j}(1-p^{(t)})^{1-y_j} + (1-\pi^{(t)}) (q^{(t)})^{y_j}(1-q^{(t)})^{1-y_j}}
\]

This $\gamma_j$ represents the probability that coin B was used for observation $j$.

The Q-function becomes:
\[
Q(\theta, \theta^{(t)}) = \sum_{j=1}^n \left[\gamma_j^{(t)} \log(\pi p^{y_j}(1-p)^{1-y_j}) + (1-\gamma_j^{(t)}) \log((1-\pi) q^{y_j}(1-q)^{1-y_j})\right]
\]

\subsection{M-step: Closed-Form Updates}

We maximize $Q(\theta, \theta^{(t)})$ by taking derivatives:

\begin{align*}
\frac{\partial Q}{\partial \pi} &= \sum_{j=1}^n \left[\frac{\gamma_j^{(t)}}{\pi} - \frac{1-\gamma_j^{(t)}}{1-\pi}\right] = 0 \\
&\Rightarrow \pi^{(t+1)} = \frac{1}{n} \sum_{j=1}^n \gamma_j^{(t)}
\end{align*}

\begin{align*}
\frac{\partial Q}{\partial p} &= \sum_{j=1}^n \gamma_j^{(t)} \left[\frac{y_j}{p} - \frac{1-y_j}{1-p}\right] = 0 \\
&\Rightarrow p^{(t+1)} = \frac{\sum_{j=1}^n \gamma_j^{(t)} y_j}{\sum_{j=1}^n \gamma_j^{(t)}}
\end{align*}

\begin{align*}
\frac{\partial Q}{\partial q} &= \sum_{j=1}^n (1-\gamma_j^{(t)}) \left[\frac{y_j}{q} - \frac{1-y_j}{1-q}\right] = 0 \\
&\Rightarrow q^{(t+1)} = \frac{\sum_{j=1}^n (1-\gamma_j^{(t)}) y_j}{\sum_{j=1}^n (1-\gamma_j^{(t)})}
\end{align*}

\subsection{Why EM Guarantees Improvement}

Let's prove why the log-likelihood never decreases. For every $z$,
$\log P(y|\theta) = \log P(y,z|\theta) - \log P(z|y,\theta)$. Taking the expectation of both sides
under $P(z|y,\theta^{(t)})$ (the left-hand side does not depend on $z$) gives, for every $\theta$,
\[
\mathcal{L}(\theta) = Q(\theta, \theta^{(t)}) + H(\theta, \theta^{(t)}).
\]
Hence
\begin{align*}
\mathcal{L}(\theta^{(t+1)}) - \mathcal{L}(\theta^{(t)})
&= [Q(\theta^{(t+1)}, \theta^{(t)}) - Q(\theta^{(t)}, \theta^{(t)})] \\
&\quad + [H(\theta^{(t+1)}, \theta^{(t)}) - H(\theta^{(t)}, \theta^{(t)})]
\end{align*}

where $H(\theta, \theta^{(t)}) = -\mathbb{E}_{P(z|y,\theta^{(t)})}[\log P(z|y,\theta)]$.

Now let's examine the second term more carefully. First, recall that at iteration $t$, we use $q(z) = P(z|y,\theta^{(t)})$. The entropy term is:

\[
H(\theta, \theta^{(t)}) = -\mathbb{E}_{P(z|y,\theta^{(t)})}[\log P(z|y,\theta)]
\]

Now consider the difference:

\begin{align*}
H(\theta^{(t+1)}, \theta^{(t)}) - H(\theta^{(t)}, \theta^{(t)}) 
&= -\mathbb{E}_{P(z|y,\theta^{(t)})}[\log P(z|y,\theta^{(t+1)})] + \mathbb{E}_{P(z|y,\theta^{(t)})}[\log P(z|y,\theta^{(t)})] \\
&= \mathbb{E}_{P(z|y,\theta^{(t)})}\left[\log\frac{P(z|y,\theta^{(t)})}{P(z|y,\theta^{(t+1)})}\right]
\end{align*}

But this is exactly the \textbf{KL divergence}:

\[
\text{KL}(P(z|y,\theta^{(t)}) \| P(z|y,\theta^{(t+1)})) = \mathbb{E}_{P(z|y,\theta^{(t)})}\left[\log\frac{P(z|y,\theta^{(t)})}{P(z|y,\theta^{(t+1)})}\right]
\]

And we know that KL divergence is always non-negative:
\[
\text{KL}(P \| Q) \geq 0 \quad \text{for any distributions } P, Q
\]
with equality if and only if $P = Q$.

Therefore:
\[
H(\theta^{(t+1)}, \theta^{(t)}) - H(\theta^{(t)}, \theta^{(t)}) = \text{KL}(P(z|y,\theta^{(t)}) \| P(z|y,\theta^{(t+1)})) \geq 0
\]

Now going back to our original difference:

\begin{align*}
\mathcal{L}(\theta^{(t+1)}) - \mathcal{L}(\theta^{(t)}) 
&= [Q(\theta^{(t+1)}, \theta^{(t)}) - Q(\theta^{(t)}, \theta^{(t)})] \\
&\quad + \underbrace{[H(\theta^{(t+1)}, \theta^{(t)}) - H(\theta^{(t)}, \theta^{(t)})]}_{\text{KL divergence} \geq 0}
\end{align*}

Since we choose $\theta^{(t+1)}$ to maximize $Q(\theta, \theta^{(t)})$ in the M-step:
\[
Q(\theta^{(t+1)}, \theta^{(t)}) - Q(\theta^{(t)}, \theta^{(t)}) \geq 0
\]

Therefore:
\[
\mathcal{L}(\theta^{(t+1)}) - \mathcal{L}(\theta^{(t)}) \geq 0
\]

\textbf{Intuition}: The log-likelihood increases by at least the improvement of the Q-function: the second term, the KL divergence between the old and new posteriors, can only add to the gain. Since the M-step does not decrease $Q$, the likelihood cannot decrease. (It is even enough to \emph{increase} $Q$ rather than maximize it: this is the ``generalized EM'' algorithm.)

\section{Python Implementation}

\subsection{Complete Implementation with Explanations}

\begin{lstlisting}[language=Python, caption=Complete Three-Coin EM Implementation]
import numpy as np
import matplotlib.pyplot as plt

class ThreeCoinEM:
    def __init__(self, n_restarts=20, tol=1e-8):
        self.n_restarts = n_restarts
        self.tol = tol
        
    def e_step(self, Y, pi, p, q):
        """
        E-step: Compute responsibilities gamma_j = P(z_j=1|y_j,theta)
        
        For each observation y_j, we compute the probability that 
        coin B was used (z_j=1) versus coin C was used (z_j=0)
        """
        gamma = np.zeros(len(Y))
        for j, y_j in enumerate(Y):
            # Probability if coin B was used
            prob_if_B = pi * (p ** y_j) * ((1 - p) ** (1 - y_j))
            # Probability if coin C was used
            prob_if_C = (1 - pi) * (q ** y_j) * ((1 - q) ** (1 - y_j))
            # Bayes' theorem: P(B|data) = P(data|B)P(B) / P(data)
            gamma[j] = prob_if_B / (prob_if_B + prob_if_C)
        return gamma
    
    def m_step(self, Y, gamma):
        """
        M-step: Update parameters to maximize Q(theta, theta_old)
        
        These updates have intuitive interpretations:
        - pi: average of responsibilities (how often coin B is selected)
        - p: success rate when coin B is used
        - q: success rate when coin C is used
        """
        n = len(Y)
        
        # pi = probability of selecting coin B
        pi_new = np.mean(gamma)
        
        # p = success rate for coin B (weighted by responsibilities)
        sum_gamma = np.sum(gamma)
        p_new = np.sum(gamma * Y) / sum_gamma if sum_gamma > 0 else 0.5
        
        # q = success rate for coin C (weighted by 1 - responsibilities)
        sum_one_minus_gamma = np.sum(1 - gamma)
        q_new = np.sum((1 - gamma) * Y) / sum_one_minus_gamma if sum_one_minus_gamma > 0 else 0.5
        
        return pi_new, p_new, q_new
    
    def log_likelihood(self, Y, pi, p, q):
        """Compute the true log-likelihood L(theta) = log P(Y|theta)"""
        log_lik = 0
        for y_j in Y:
            prob_B = pi * (p ** y_j) * ((1 - p) ** (1 - y_j))
            prob_C = (1 - pi) * (q ** y_j) * ((1 - q) ** (1 - y_j))
            log_lik += np.log(prob_B + prob_C)
        return log_lik
    
    def Q_function(self, Y, gamma, pi, p, q):
        """
        Compute Q(theta, theta_old) = E[log P(Y,Z|theta)]
        
        This is the expected complete data log-likelihood,
        where expectation is taken over Z given current parameters
        """
        Q_val = 0
        for j, y_j in enumerate(Y):
            # Contribution if coin B was used
            term_B = gamma[j] * np.log(pi * (p ** y_j) * ((1 - p) ** (1 - y_j)))
            # Contribution if coin C was used
            term_C = (1 - gamma[j]) * np.log((1 - pi) * (q ** y_j) * ((1 - q) ** (1 - y_j)))
            Q_val += term_B + term_C
        return Q_val
    
    def fit(self, Y, verbose=True):
        """
        Run EM algorithm with multiple random restarts
        
        Why multiple restarts? EM converges to local optima,
        so we try many starting points and pick the best one.
        """
        best_solution = None
        best_ll = -np.inf
        
        if verbose:
            print("Running EM with multiple initializations:")
            print("Run |     π    |     p    |     q    | Log-Likelihood")
            print("----|----------|----------|----------|---------------")
        
        for run in range(self.n_restarts):
            # Random initialization (avoid boundaries)
            pi, p, q = np.random.uniform(0.1, 0.9, 3)
            
            # Track progress for analysis
            history = {
                'pi': [pi], 'p': [p], 'q': [q], 
                'Q': [], 'L': [], 'iterations': [],
                'gamma_history': []
            }
            
            # EM iterations
            for iteration in range(100):
                # E-step: Compute responsibilities
                gamma = self.e_step(Y, pi, p, q)
                history['gamma_history'].append(gamma.copy())
                
                # Compute Q and L for monitoring
                Q_val = self.Q_function(Y, gamma, pi, p, q)
                L_val = self.log_likelihood(Y, pi, p, q)
                history['Q'].append(Q_val)
                history['L'].append(L_val)
                history['iterations'].append(iteration)
                
                # M-step: Update parameters
                pi_new, p_new, q_new = self.m_step(Y, gamma)
                
                # Check convergence
                param_change = abs(pi_new - pi) + abs(p_new - p) + abs(q_new - q)
                if param_change < self.tol:
                    break
                    
                pi, p, q = pi_new, p_new, q_new
                history['pi'].append(pi)
                history['p'].append(p)
                history['q'].append(q)
            
            final_ll = self.log_likelihood(Y, pi, p, q)
            if final_ll > best_ll:
                best_ll = final_ll
                best_solution = (pi, p, q, history)
                
            if verbose:
                print(f"{run:3d} | {pi:.6f} | {p:.6f} | {q:.6f} | {final_ll:.6f}")
        
        return best_solution

# Let's run it on our original data!
Y = np.array([1, 1, 0, 1, 0, 0, 1, 0, 1, 1])
print("Original data:", Y)
print(f"Number of heads: {np.sum(Y)}/{len(Y)}")

model = ThreeCoinEM(n_restarts=10)
pi, p, q, history = model.fit(Y)

print(f"\n=== FINAL RESULTS ===")
print(f"π (prob of selecting coin B) = {pi:.4f}")
print(f"p (prob of heads for coin B) = {p:.4f}") 
print(f"q (prob of heads for coin C) = {q:.4f}")
print(f"Final log-likelihood = {model.log_likelihood(Y, pi, p, q):.4f}")
\end{lstlisting}

\subsection{Visualization and Analysis}

\begin{lstlisting}[language=Python, caption=Visualizing EM Progress]
def plot_em_progress(history, Y):
    """Plot the EM algorithm's progress"""
    fig, ((ax1, ax2), (ax3, ax4)) = plt.subplots(2, 2, figsize=(12, 8))
    
    # Plot 1: Parameter convergence
    iterations = range(len(history['pi']))
    ax1.plot(iterations, history['pi'], 'b-', label='π', linewidth=2, marker='o', markersize=4)
    ax1.plot(iterations, history['p'], 'r-', label='p', linewidth=2, marker='s', markersize=4)
    ax1.plot(iterations, history['q'], 'g-', label='q', linewidth=2, marker='^', markersize=4)
    ax1.set_xlabel('Iteration')
    ax1.set_ylabel('Parameter Value')
    ax1.set_title('Parameter Convergence')
    ax1.legend()
    ax1.grid(True, alpha=0.3)
    
    # Plot 2: Q-function and true likelihood
    ax2.plot(history['iterations'], history['Q'], 'orange', 
             label='Q(θ,θ^(t))', linewidth=2)
    ax2.plot(history['iterations'], history['L'], 'purple', 
             label='L(θ) = log P(Y|θ)', linewidth=2)
    ax2.set_xlabel('Iteration')
    ax2.set_ylabel('Log-Likelihood')
    ax2.set_title('Q-Function vs True Likelihood')
    ax2.legend()
    ax2.grid(True, alpha=0.3)
    
    # Plot 3: Bound tightness
    gap = [L - Q for L, Q in zip(history['L'], history['Q'])]
    ax3.plot(history['iterations'], gap, 'red', linewidth=2)
    ax3.set_xlabel('Iteration')
    ax3.set_ylabel('L(θ) - Q(θ,θ^(t))')
    ax3.set_title('Tightness of Lower Bound (KL Divergence)')
    ax3.grid(True, alpha=0.3)
    
    # Plot 4: Final responsibilities
    final_gamma = history['gamma_history'][-1]
    colors = ['red' if y == 0 else 'blue' for y in Y]
    ax4.bar(range(len(Y)), final_gamma, color=colors, alpha=0.7)
    ax4.set_xlabel('Observation Index')
    ax4.set_ylabel('P(coin B was used)')
    ax4.set_title('Final Responsibilities\n(Blue = heads, Red = tails)')
    ax4.grid(True, alpha=0.3)
    
    plt.tight_layout()
    plt.show()

# Generate some analysis
print("\n=== INTERPRETATION ===")
print(f"According to our model:")
print(f"- Coin A selects coin B with probability {pi:.2f} and coin C with probability {1-pi:.2f}")
print(f"- Coin B shows heads with probability {p:.2f}")
print(f"- Coin C shows heads with probability {q:.2f}")
print(f"- Overall expected heads: {pi*p + (1-pi)*q:.2f}")
print(f"- Actual observed heads: {np.mean(Y):.2f}")

plot_em_progress(history, Y)
\end{lstlisting}

\section{Deep Dive: Understanding EM and ELBO}

\subsection{The Relationship Between EM and ELBO}

The EM algorithm can be viewed as \textbf{coordinate ascent} on the ELBO:

\begin{itemize}
    \item \textbf{E-step}: Maximize ELBO with respect to $q$ (distribution over latent variables)
    \[
    q^{(t+1)} = \arg\max_q \text{ELBO}(\theta^{(t)}, q) = P(z|y,\theta^{(t)})
    \]
    
    \item \textbf{M-step}: Maximize ELBO with respect to $\theta$ (model parameters)
    \[
    \theta^{(t+1)} = \arg\max_{\theta} \text{ELBO}(\theta, q^{(t+1)}) = \arg\max_{\theta} Q(\theta, \theta^{(t)})
    \]
\end{itemize}

The ELBO decomposes as:
\[
\text{ELBO}(\theta, q) = \mathbb{E}_q[\log P(y,z|\theta)] - \mathbb{E}_q[\log q(z)] = \log P(y|\theta) - \text{KL}(q(z) \| P(z|y,\theta))
\]

So maximizing ELBO is equivalent to:
\begin{enumerate}
    \item Maximizing the expected complete data likelihood
    \item Minimizing the KL divergence between $q(z)$ and the true posterior
\end{enumerate}

\subsection{Why This Matters in Modern Machine Learning}

The EM/ELBO framework is fundamental to many modern techniques:
\begin{itemize}
    \item \textbf{Variational Autoencoders (VAEs)}: Use the ELBO to train deep generative models
    \item \textbf{Variational Inference}: General framework for approximate Bayesian inference
    \item \textbf{Gaussian Mixture Models}: EM is the standard training algorithm
    \item \textbf{Hidden Markov Models}: Use EM (Baum-Welch algorithm) for parameter estimation
\end{itemize}

\section{Common Questions and Pitfalls}

\subsection{Why Multiple Restarts?}

EM is guaranteed not to decrease the likelihood, and under mild regularity conditions it converges to a \textbf{stationary point} of the likelihood (usually a local maximum, possibly a saddle point), not necessarily the global maximum. The likelihood surface can have multiple peaks, so we try different starting points.

\subsection{Identifiability Issue}

Notice that if we swap $p$ and $q$ and change $\pi$ to $1-\pi$, we get the same likelihood:
\[
\mathcal{L}(\pi, p, q) = \mathcal{L}(1-\pi, q, p)
\]

This means there are (at least) two equivalent solutions! In fact the situation is much more
degenerate for the three-coin model: each observation is a single Bernoulli variable with
\[
P(y = 1 \mid \theta) = \pi p + (1 - \pi) q =: s(\theta),
\]
so the likelihood $\mathcal{L}(\theta) = n_1 \log s(\theta) + n_0 \log(1 - s(\theta))$ depends on the three
parameters only through $s(\theta)$. It is maximized by \emph{every} $\theta$ with $s(\theta) = n_1/n$
($= 0.6$ for our data): a whole two-dimensional surface of maximum likelihood estimates. The model is not
identifiable from these data, and EM simply converges to one point of this surface, determined by the
initialization. (Identifiability would require more information, e.g.\ several tosses of the selected
coin for each toss of coin A.)

\subsection{Sensitivity to Initialization}

Different initializations lead to very different parameter estimates, all with the same likelihood. Here this is not a small-sample effect but a consequence of the non-identifiability explained above; with more data the estimates would \emph{not} become unique.

\section{Conclusion}

The three-coin problem teaches us several important lessons:

\begin{itemize}
    \item \textbf{EM provides a principled approach} to maximum likelihood estimation with latent variables
    \item \textbf{The ELBO framework} gives us a systematic way to handle incomplete data
        \item \textbf{Monotonic improvement} of the likelihood is guaranteed, but convergence is only to a stationary point (often a local optimum)
    \item \textbf{Multiple restarts} are essential in practice
    \item \textbf{These ideas generalize} to much more complex models
\end{itemize}

The beauty of EM is that it turns a hard optimization problem into an iterative procedure where each step is easy. This pattern appears throughout machine learning!
\end{notechapter}
